\section{Experimental Setup}\label{sec:experiments}
Describe how you set up your experiments to validate the proposed model.

\subsection{Datasets}
Describe the datasets used in your experiments. Provide key details (e.g., number of samples, features, classes, train/val/test splits) in a summary table.

\begin{table}[h]
\caption{Summary of experimental datasets.}\label{tab:datasets}
\begin{tabular}{@{}lcccc@{}}
\toprule
Dataset Name & Number of Samples & Features & Classes & Train/Val/Test Split \\
\midrule
Dataset 1 & 10,000 & 128 & 10 & 70\% / 10\% / 20\% \\
Dataset 2 & 50,000 & 256 & 2 & 80\% / 10\% / 10\% \\
Dataset 3 & 150,000 & 512 & 100 & 60\% / 20\% / 20\% \\
\bottomrule
\end{tabular}
\end{table}

\subsection{Evaluation Metrics}
Define the metrics used to evaluate the model performance (e.g., Accuracy, F1-Score, Mean Squared Error, Area Under Curve). Write equations where helpful.
\begin{equation}
\text{F1-Score} = 2 \cdot \frac{\text{Precision} \cdot \text{Recall}}{\text{Precision} + \text{Recall}}
\label{eq:f1_score}
\end{equation}

\subsection{Baseline Models}
Briefly explain the baseline models you compare your approach against:
\begin{itemize}
    \item \textbf{Baseline 1 (e.g., Random Forest / Support Vector Machine)}: Traditional ML models.
    \item \textbf{Baseline 2 (e.g., standard CNN / LSTM)}: Standard deep learning configurations.
    \item \textbf{Baseline 3 (e.g., State-of-the-art transformer-based method)}: The current leading model in this field.
\end{itemize}

\subsection{Implementation Details}
Provide details about your hardware, software framework, hyperparameter tuning, optimizer, and training parameters (to ensure reproducibility):
\begin{itemize}
    \item \textbf{Framework}: PyTorch 2.x / TensorFlow 2.x
    \item \textbf{Hardware}: NVIDIA RTX 4090 GPU / A100 GPU
    \item \textbf{Optimizer}: AdamW with learning rate $\eta = 10^{-4}$ and weight decay $10^{-2}$.
    \item \textbf{Hyperparameters}: Batch size $B = 64$, dropout rate $p = 0.2$.
\end{itemize}
